\documentclass[a4paper]{article}
% Español (México) con XeLaTeX: fontspec para fuentes Unicode, polyglossia
% para la división silábica y los nombres del documento en la variante mexicana.
\usepackage{fontspec}
\usepackage{polyglossia}
\setdefaultlanguage[variant=mexican]{spanish}
% Los nombres chinos van como «pinyin (original)»: xeCJK da a los caracteres
% Han su propia fuente; sin ella XeLaTeX los omite con un aviso.
% FandolSong no cubre algunos caracteres de nombres (U+73FA); WenQuanYi Zen Hei sí.
\usepackage[AutoFallBack=true]{xeCJK}
\setCJKmainfont{FandolSong-Regular.otf}
\setCJKfallbackfamilyfont{\CJKrmdefault}{WenQuanYi Zen Hei}
% polyglossia no traduce los nombres de listings.
\AtBeginDocument{\providecommand{\lstlistingname}{}\renewcommand{\lstlistingname}{Listado}%
  \providecommand{\lstlistlistingname}{}\renewcommand{\lstlistlistingname}{Índice de listados}}
\usepackage{amsmath, amssymb}
\usepackage{graphicx}
\usepackage[margin=2.5cm]{geometry}
\usepackage[most]{tcolorbox}
\usepackage{etoolbox}
\usepackage{listings}
\usepackage{hyperref}
\usepackage{booktabs}
\usepackage{float}
\newtcolorbox{knowledgebox}[1]{enhanced, colback=blue!5!white, colframe=blue!75!black, colbacktitle=blue!75!black, coltitle=white, fonttitle=\bfseries, title={#1}, attach boxed title to top left={yshift=-2mm, xshift=2mm}, boxrule=1pt, sharp corners}
\newtcolorbox{importantbox}[1]{enhanced, colback=yellow!10!white, colframe=yellow!80!black, colbacktitle=yellow!80!black, coltitle=black, fonttitle=\bfseries, title={#1}, sharp corners}
\newtcolorbox{warningbox}[1]{enhanced, colback=red!5!white, colframe=red!75!black, colbacktitle=red!75!black, coltitle=white, fonttitle=\bfseries, title={#1}, sharp corners}
\lstset{language=Python, basicstyle=\ttfamily\small, keywordstyle=\color{blue}, stringstyle=\color{red!60!black}, commentstyle=\color{green!60!black}, breaklines=true, frame=single, numbers=left, numberstyle=\tiny\color{gray}, captionpos=b, extendedchars=true}
\newcommand{\notetitle}{Inferencia y despliegue de vLLM: ajuste de parámetros y análisis de memoria de video}
\newcommand{\noteauthors}{Elaboradas a partir de materiales públicos del curso}
\newcommand{\notedate}{2025}
\newcommand{\videochannel}{Wudaokou Nash (五道口纳什)}
\newcommand{\videopublishdate}{2025}
\newcommand{\videoduration}{aproximadamente 20 minutos}
\newcommand{\videourl}{}
\newcommand{\videocoverpath}{cover.jpg}
\newcommand{\repourl}{https://github.com/hqhq1025/ai-course-notes}

% Footer with repo info
\usepackage{fancyhdr}
\pagestyle{fancy}
\fancyhf{}
\fancyfoot[C]{\thepage}
\fancyfoot[R]{\footnotesize\color{gray}\href{\repourl}{AI Course Notes}}
\renewcommand{\headrulewidth}{0pt}
\renewcommand{\footrulewidth}{0pt}

\begin{document}
\begin{titlepage}\centering
{\Large Notas del curso\par}\vspace{1.2cm}{\huge\bfseries \notetitle\par}\vspace{0.8cm}{\large \noteauthors\par}\vspace{0.3cm}{\large \notedate\par}\vspace{1.2cm}
\ifdefempty{\videocoverpath}{}{\includegraphics[width=0.82\textwidth,height=0.45\textheight,keepaspectratio]{\videocoverpath}\par}
\vfill
\begin{tcolorbox}[width=0.9\textwidth, colback=black!2!white, colframe=black!60, sharp corners]
\textbf{Autor o canal del video}: \videochannel\par\textbf{Fecha de publicación}: \videopublishdate\par\textbf{Duración del video}: \videoduration\par
\ifdefempty{\videourl}{}{\textbf{Enlace del video}: \href{\videourl}{\nolinkurl{\videourl}}\par}\par
\textbf{Página del proyecto}: \href{\repourl}{\nolinkurl{\repourl}}\par
\end{tcolorbox}
\end{titlepage}
\tableofcontents\newpage

\section{Introducción}
En esta clase profundizamos en más parámetros y detalles de la inferencia y el despliegue de vLLM, incluidos el análisis de la memoria de GPU y los métodos de ajuste. vLLM tiene dos grandes escenarios de aplicación: el despliegue de inferencia y el rollout en línea dentro de RL Training.

\begin{importantbox}{Los dos grandes escenarios de vLLM}
\begin{itemize}
  \item \textbf{Despliegue de inferencia}: Offline Batch Generation + Online Serving (compatible con la API de OpenAI)
  \item \textbf{RL Training}: rollout en línea del modelo dentro de RL for Language Model (por ejemplo, OpenRLHF fue el primero en introducir un motor de inferencia moderno en el pipeline de RL)
\end{itemize}
\end{importantbox}

\section{Online vs. Offline 推理}
\subsection{Offline Batch Generation}
Adecuado para escenarios de procesamiento por lotes, que procesan una gran cantidad de solicitudes de una sola vez para maximizar el rendimiento.

\subsection{Online Serving}
Se despliega como un servicio compatible con la API de OpenAI, que admite solicitudes en tiempo real. Métricas clave: TTFT (latencia hasta el primer token) y TPS (tokens por segundo).

\section{Análisis de la memoria de la GPU}
Pesos del modelo + KV Cache + valores de activación intermedios. El KV Cache crece linealmente con la longitud de la secuencia y es el principal cuello de botella de la inferencia con contexto largo.

\subsection{Resumen de la sección}
El consumo de memoria de la GPU en vLLM depende principalmente de los pesos del modelo y del KV Cache. Una configuración adecuada de los parámetros es fundamental para una inferencia eficiente.

\section{Perspectiva de sistema: qué revisar primero cuando el rendimiento está limitado}
Para el entrenamiento de Agentic RL, el modo de inferencia y el uso de memoria de video no son dos temas independientes. El verdadero cuello de botella del sistema suele manifestarse así: por un lado se quiere aumentar el rendimiento del rollout, y por otro el context, el batch y la cache saturan la memoria de video. Entender la diferencia entre online y offline es, en esencia, entender cómo el sistema de entrenamiento negocia entre latencia y rendimiento.

\begin{knowledgebox}{Orden de revisión sugerido}
\begin{itemize}
  \item Primero revisa el número de tokens y el nivel de concurrencia de cada ronda de rollout
  \item Después revisa si la KV cache y el tamaño de lote saturan la memoria de video
  \item Por último decide si usar una orquestación de inferencia online u offline
\end{itemize}
\end{knowledgebox}

\subsection{Resumen de la sección}
El análisis de memoria de video no es un tema secundario, sino parte del diseño del sistema de RL. Solo cuando el rendimiento, la latencia y la memoria de video están equilibrados a la vez, la canalización de entrenamiento puede funcionar de manera estable.

\newpage
\section{Síntesis y ampliación}
\begin{enumerate}
  \item vLLM presta servicio tanto al despliegue de inferencia como al RL Training
  \item Los modos Offline y Online tienen cada uno su propia dirección de optimización
  \item Análisis de memoria de video: pesos del modelo + KV Cache + valores de activación
\end{enumerate}
\end{document}
